\documentclass[11pt,a4paper]{article}
\usepackage[T2A]{fontenc}
\usepackage[utf8]{inputenc}
\usepackage[russian,english]{babel}
\usepackage{amsmath,amssymb,amsthm,mathtools}
\usepackage{setspace}
\usepackage{enumitem}
\usepackage[hidelinks]{hyperref}
\onehalfspacing
\newcommand{\head}[1]{\par\medskip\noindent\textbf{\underline{#1}}\quad}
\newcommand{\sheettitle}[3]{% title, subtitle, homework-flag (empty or "Homework")
  \begin{center}
  {\huge\bfseries #1\par}\vspace{1.2em}
  {\Large #2\par}\vspace{0.8em}
  \if\relax\detokenize{#3}\relax\else{\Large #3\par}\vspace{0.8em}\fi
  {\foreignlanguage{russian}{\rudate}\par}
  \end{center}\vspace{0.5em}}
\newcommand{\src}[1]{\hfill(#1)}
\newcommand{\rudate}{\number\day~\ifcase\month\or января\or февраля\or марта\or апреля\or мая\or июня\or июля\or августа\or сентября\or октября\or ноября\or декабря\fi~\number\year~года}
\begin{document}
\sheettitle{Polynomials -- 9}{Cyclotomic polynomials and the arithmetic of roots of unity}{Homework}

\head{Theoretical minimum.} Cyclotomic polynomials $\Phi_n$: integer coefficients,
$x^n-1=\prod_{d\mid n}\Phi_d$, the M\"obius formula $\Phi_n=\prod_{d\mid n}(x^d-1)^{\mu(n/d)}$,
the values $\Phi_n(1)$, the formulas for $\Phi_p$ and $\Phi_{pq}$; irreducibility over
$\mathbb{Q}$, so $\Phi_n$ is the minimal polynomial of a primitive $n$-th root of unity $\zeta$
and $\sum a_k\zeta^k=0$ ($a_k\in\mathbb{Q}$) iff $\Phi_n\mid\sum a_kx^k$; non-negative vanishing
sums of $pq$-th roots are unions of regular $p$- and $q$-gons, and the Lam--Leung theorem. Over
$\mathbb{F}_p$: the power sums $\sum_x x^k$, polynomial functions of degree $<p$, Hermite's
criterion for permutation polynomials. Annihilating ideals of the shift operator
$f(x)\mapsto f(x+h)$ are principal and closed under gcd.

\medskip\noindent\rule{\linewidth}{0.4pt}

\begin{enumerate}[leftmargin=2.2em,itemsep=0.6em]
\item Prove the formula for $\sum_{x\in\mathbb{F}_p}x^k$ and Hermite's criterion. Deduce that
$x^k$ permutes $\mathbb{F}_p$ if and only if $\gcd(k,p-1)=1$, and that no polynomial of degree
$d>1$ with $d\mid p-1$ is a permutation polynomial of $\mathbb{F}_p$.

\item Let $p$ be a prime. Prove that $\Phi_{pn}(x)=\Phi_n(x^p)$ if $p\mid n$, that
$\Phi_{pn}(x)=\Phi_n(x^p)/\Phi_n(x)$ if $p\nmid n$, and that $\Phi_{2n}(x)=\Phi_n(-x)$ for odd
$n>1$. For distinct primes $p,q$ prove that all coefficients of $\Phi_{pq}$ lie in
$\{-1,0,1\}$ by expanding $\frac{1-x}{(1-x^p)(1-x^q)}$ modulo $x^{pq}$, and show that the coefficient of $x^7$
in $\Phi_{105}$ equals $-2$.

\item (Kronecker) Let $f\in\mathbb{Z}[x]$ be monic with $f(0)\neq0$, all of whose complex roots
satisfy $|z|\le1$. Prove that all roots of $f$ are roots of unity. (Consider the polynomials
$f_m=\prod(x-z_i^m)$.) Deduce that an algebraic integer all of whose conjugates have absolute
value $1$ is a root of unity, and show by the example $\frac{3+4i}{5}$ that integrality cannot be
dropped.

\item Find all pairs of cubical dice whose faces are labelled by positive integers (repetitions
allowed) such that the sum of the two numbers thrown has the same distribution as for two
ordinary dice. (Factor the generating function $(x+x^2+\dots+x^6)^2$ into cyclotomic
polynomials.)

\item A family of at least two sets is called a partition of a set $X$ if the sets are pairwise
disjoint and their union is $X$.
\begin{enumerate}[label=\alph*)]
\item Construct a partition of the set of all positive integers into infinitely many disjoint
infinite arithmetic progressions with different common differences.
\item Can the set of all positive integers be partitioned into a finite number of such
progressions?
\end{enumerate}
\src{ISTCiM 2022}

\item Determine all pairs $P(x),Q(x)$ of complex polynomials with leading coefficient $1$ such
that $P(x)$ divides $Q(x)^2+1$ and $Q(x)$ divides $P(x)^2+1$.\\ \mbox{}\src{IMC 2018}

\item Let $m$ be a positive integer and let $p$ be a prime divisor of $m$. Suppose that the
complex polynomial $a_0+a_1x+\dots+a_nx^n$ with $n<\frac{p}{p-1}\varphi(m)$ and $a_n\neq0$ is
divisible by the cyclotomic polynomial $\Phi_m(x)$. Prove that there are at least $p$ nonzero
coefficients $a_i$. \src{VJIMC 2015}

\item Suppose that for a function $f\colon\mathbb{R}\to\mathbb{R}$ and real numbers $a<b$ one
has $f(x)=0$ for all $x\in(a,b)$. Prove that $f(x)=0$ for all $x\in\mathbb{R}$ if
\[
\sum_{k=0}^{p-1}f\Bigl(y+\frac kp\Bigr)=0
\]
for every prime number $p$ and every real number $y$. \src{IMC 2010}

\item Let $p$ be an odd prime. Show that for at least $(p+1)/2$ values of $n$ in
$\{0,1,2,\dots,p-1\}$ the number
\[
\sum_{k=0}^{p-1}k!\,n^k
\]
is not divisible by $p$ (here $0^0=1$). \src{Putnam 2011 B6}

\item Let $k$ and $n$ be positive integers with $n\ge k^2-3k+4$, and let
\[
f(z)=z^{n-1}+c_{n-2}z^{n-2}+\dots+c_0
\]
be a polynomial with complex coefficients such that
$c_0c_{n-2}=c_1c_{n-3}=\dots=c_{n-2}c_0=0$. Prove that $f(z)$ and $z^n-1$ have at most $n-k$
common roots. \src{IMC 2017}
\end{enumerate}
\end{document}
